\section*{Chapitre \ref{ch:ll:memory-hierarchy-and-caches} --- Hiérarchie de la mémoire et caches}

\begin{solution}{exo:ll:memory-hierarchy-and-caches:1}
Le décalage est formé des 6 bits de poids faible~: $\texttt{0x48} \bmod 64 = 8$. L'ensemble est formé des bits 6 à 11~: $\lfloor\texttt{0x5e48}/64\rfloor
\bmod 64 = 377 \bmod 64 = 57$.
\end{solution}

\begin{solution}{exo:ll:memory-hierarchy-and-caches:2}
$2\,097\,152 / (64 \times 16) = 2\,048$ ensembles~; des adresses distantes de $2\,048 \times 64 = 131\,072$ octets (128~KiB) partagent le même ensemble.
\end{solution}

\begin{solution}{exo:ll:memory-hierarchy-and-caches:3}
$600~\text{MiB}/4~\text{KiB} = 153\,600$ entrées avec des pages de base, bien au-delà de la capacité de n'importe quel TLB~; 300 avec des pages de 2~MiB.
\end{solution}

\begin{solution}{exo:ll:memory-hierarchy-and-caches:4}
\omterm{def:ll:memory-hierarchy-and-caches:false}{Faux partage} (false sharing)~: chaque écriture d'un thread invalide la copie de la ligne détenue par l'autre, et les deux paient un transfert entre cœurs
à la plupart des écritures (environ cinq fois le coût d'un incrément atomique sans contention dans la
\cref{fig:ll:memory-hierarchy-and-caches:false}). Placer chaque champ sur sa propre \omterm{def:ll:memory-hierarchy-and-caches:line}{ligne de cache}~: 128 octets au lieu de 16.
\end{solution}

\begin{solution}{exo:ll:memory-hierarchy-and-caches:5}
Environ $4 \times 3.5 = \qty{14}{\nano\second}$ depuis le second niveau et $4 \times 150 = \qty{600}{\nano\second}$ depuis
la mémoire~: les chargements dépendants ne peuvent pas se chevaucher.
\end{solution}

\begin{solution}{exo:ll:memory-hierarchy-and-caches:6}
Dans l'ordre des adresses, la ligne suivante est prévisible~: le \omterm{def:ll:memory-hierarchy-and-caches:prefetch}{préchargement matériel} la charge avant qu'elle ne soit nécessaire, et les chargements se chevauchent~;
dans un ordre aléatoire, chaque adresse n'est connue qu'au retour du chargement précédent. Les structures chaînées par pointeurs (listes, arbres,
tables associatives à nœuds) placent leurs nœuds là où l'allocateur les place, si bien que les parcours sont des suites d'accès aléatoires~: sur le chemin critique (hot path), il faut privilégier
des tableaux contigus.
\end{solution}

\begin{solution}{exo:ll:memory-hierarchy-and-caches:7}
La vérification des bornes est une comparaison et un branchement toujours prédit non pris, indépendant du résultat du chargement, de sorte que
le cœur l'exécute pendant que le chargement dépendant qu'il protège est en attente. Sur un anneau de 64~MiB, chaque pas attend la mémoire
(environ \qty{150}{\nano\second} ici)~; une vérification d'un ou deux cycles ne se mesure pas. Elle se mesurerait sur un anneau qui tient dans le premier
niveau, où le chargement lui-même prend cinq cycles.
\end{solution}

\begin{solution}{exo:ll:memory-hierarchy-and-caches:8}
Un seul instrument et les mêmes quelques messages gardent tout l'\omterm{def:ll:memory-hierarchy-and-caches:miss}{ensemble de travail} (working set) dans le premier niveau et le prédicteur entraîné~: le
\qty{20}{\nano\second} est la valeur du premier niveau. En production, les mises à jour touchent de nombreux instruments et de nombreux niveaux de prix,
dont la plupart ne tiennent pas dans le premier niveau et certains pas dans le second. Rejouer une journée enregistrée sur tous les instruments et
en rapporter les percentiles.
\end{solution}

\begin{solution}{pb:ll:memory-hierarchy-and-caches:1}
\begin{enumerate}
\item Environ \qty{1.1}{\nano\second}, \qty{3.5}{\nano\second} et de 100 à \qty{175}{\nano\second}.
\item Environ 48~KiB, un mégaoctet et demi à deux mégaoctets, et entre huit et douze mégaoctets~: un premier niveau de 48~KiB,
un second niveau d'environ 2~MiB, et un troisième niveau que le parcours utilise jusqu'à environ huit mégaoctets.
\item Les 24~MiB sont partagés par tous les cœurs et contiennent aussi leurs données, et un parcours aléatoire à travers les pages de la machine virtuelle
répartit les lignes sur les ensembles de façon inégale, si bien que certains ensembles débordent bien avant que tout le cache soit plein.
\item La latence du troisième niveau, 75 cycles ou plus (environ \qty{20}{\nano\second})~; mesurée, de 15 à
\qty{23}{\nano\second}.
\item $2 \times 256 \times 32 + 4\,096 = 20\,480$ octets.
\item 320 lignes.
\item $0.75 \times 2\,097\,152 / 20\,480 = 76.8$~: 76 instruments.
\item Deux au plus, sans rien d'autre dans le premier niveau.
\item $1 - 76/300 = 74.7~\%$.
\item $3 \times (0.253 \times 3.5 + 0.747 \times 150) \approx \qty{339}{\nano\second}$.
\item Les 300 carnets d'ordres occupent environ 6~Mo, soit 1\,500 pages de base, davantage que ce que couvre le TLB de premier niveau~; avec des pages de 2~MiB, trois
entrées suffisent. Le gain est plus grand sur un hôte virtualisé, où chaque parcours est bidimensionnel.
\item Ne stocker que les niveaux les plus proches du meilleur cours (une échelle plus étroite, chapitre~19) et utiliser des champs plus petits (prix et
quantités sur 32~bits, exprimés en ticks et en lots)~; garder l'état froid dans une autre structure.
\item \textbf{Résultat à retenir.} Sur ce portable~: environ \qty{1.1}{\nano\second} depuis le premier niveau,
\qty{3.5}{\nano\second} depuis le second, de 100 à \qty{175}{\nano\second} depuis la mémoire~; les carnets d'ordres de 76 instruments tiennent dans
les trois quarts d'un second niveau de 2~MiB.
\item 32 divise 64, donc deux niveaux partagent une ligne et aucun ne chevauche deux lignes~; avec 48 octets, la moitié des niveaux chevauchent une
frontière de ligne et coûtent deux \omterm{def:ll:memory-hierarchy-and-caches:miss}{défauts de cache}.
\item Quatre tableaux d'un même champ (structure de tableaux, structure of arrays) lorsqu'un balayage lit un champ sur de nombreux niveaux, par exemple pour une somme de profondeur~; une
structure lorsqu'une mise à jour lit tous les champs d'un même niveau.
\item La latence de chaque accès pris isolément, sans préchargement ni chevauchement, que masque un banc d'essai séquentiel.
\item Afin que le temps par pas soit fixé par la taille du tampon et non par un cycle plus court qui tient dans un niveau plus petit.
\item Plusieurs fois le coût d'une mise à jour sans contention, soit quelques dizaines de nanosecondes à chaque fois que l'autre thread a écrit la ligne.
\item Mesurer les percentiles de latence de la mise à jour lors d'un rejeu, les comparer à la latence du second niveau et, là où des compteurs matériels
sont disponibles, compter les \omterm{def:ll:memory-hierarchy-and-caches:miss}{défauts de cache} du second niveau.
\item Où se trouvent ses octets~: combien de chargements dépendants fait une mise à jour, et depuis quel niveau chacun est servi.
\end{enumerate}
\end{solution}

\begin{solution}{iq:ll:memory-hierarchy-and-caches:1}
Un vecteur est contigu~: ses éléments partagent des \omterm{def:ll:memory-hierarchy-and-caches:line}{lignes de cache} et le \omterm{def:ll:memory-hierarchy-and-caches:prefetch}{préchargement matériel} les charge en flux. Les nœuds d'une liste sont dispersés, chacun exigeant
un chargement dépendant, et donc un \omterm{def:ll:memory-hierarchy-and-caches:miss}{défaut de cache} potentiel.
\iqlookfor{contiguïté, préchargement et chargements dépendants.}
\end{solution}

\begin{solution}{iq:ll:memory-hierarchy-and-caches:2}
Deux threads écrivent des variables différentes dans la même \omterm{def:ll:memory-hierarchy-and-caches:line}{ligne de cache}, si bien que la ligne fait des allers-retours entre leurs cœurs. Pour la détecter, il faut
consulter le rapport de contention de cache du profileur (\texttt{perf c2c}) ou mesurer le temps avec du bourrage~; la corriger en alignant
les variables écrites indépendamment sur des lignes distinctes.
\iqlookfor{le mécanisme de \omterm{def:ll:memory-hierarchy-and-caches:coherence}{cohérence de cache}, et le bourrage ou le regroupement par thread écrivain.}
\end{solution}

\begin{solution}{iq:ll:memory-hierarchy-and-caches:3}
Le \omterm{def:ll:memory-hierarchy-and-caches:tlb}{tampon de traduction d'adresses} (TLB), cache des traductions entre adresses virtuelles et physiques~; un défaut coûte un parcours de tables de pages de plusieurs chargements dépendants, davantage sous
virtualisation. Les \omterm{def:ll:memory-hierarchy-and-caches:tlb}{grandes pages} (huge pages) permettent à une seule entrée de couvrir 2~MiB ou 1~GiB, de sorte que les traductions des données chaudes restent en cache.
\iqlookfor{portée = nombre d'entrées multiplié par la taille de page, et parcours de tables de pages.}
\end{solution}

\begin{solution}{iq:ll:memory-hierarchy-and-caches:4}
Un parcours de pointeurs~: un cycle aléatoire unique à travers un tampon, un emplacement par ligne, qui mesure des chargements dépendants pour des tailles allant de quelques
kilo-octets jusqu'au-delà du dernier niveau~; le temps par pas, en fonction de la taille, forme un escalier. Épingler le thread, utiliser des \omterm{def:ll:memory-hierarchy-and-caches:tlb}{grandes pages} pour
isoler les effets du TLB, répéter.
\iqlookfor{chargements aléatoires dépendants qui déjouent le préchargement et le chevauchement.}
\end{solution}

\begin{solution}{iq:ll:memory-hierarchy-and-caches:5}
Des tableaux de 32~KiB chacun, alloués alignés sur la page, de sorte que l'élément $i$ des deux tableaux tombe dans le même ensemble de cache (4~KiB
d'aliasage)~: avec plus de flux de ce type que de voies, les lignes s'évincent mutuellement. Décaler l'un des tableaux d'une ou deux lignes.
\iqlookfor{conflits d'ensembles dus aux pas en puissances de deux.}
\end{solution}

\begin{solution}{iq:ll:memory-hierarchy-and-caches:6}
Une table de hachage à adressage ouvert (open addressing) avec sondage linéaire, clés et valeurs dans un seul tableau plat de petits emplacements (identifiant d'ordre et
index de 32~bits dans une réserve d'ordres), dimensionnée à une puissance de deux avec un faible facteur de charge, soutenue par des \omterm{def:ll:memory-hierarchy-and-caches:tlb}{grandes pages} et
pré-chargées (pre-faulting)~; un \omterm{def:ll:memory-hierarchy-and-caches:miss}{défaut de cache} attendu par recherche, aucune allocation, aucun parcours de pointeurs. Le chapitre~19 la construit.
\iqlookfor{stockage plat, un défaut par recherche, aucune allocation à l'insertion.}
\end{solution}
